\scanpage{342}
\mspage{81}
% archive sheet number: 341
qui est bel et bien fibrant (comme produit cartésien de deux foncteurs fibrants) et a les fibres qu'il faut~! On trouve donc bien, par le \unsure{sorite} général sur les foncteurs fibrants (ici sur $\underline{\mathrm{Ty}}\times\underline{\mathrm{Ty}}$) \unsure{une|un} foncteur de catégories fibrées sur $\underline{\mathrm{Ty}}\times\mathrm{Ty}$
\diagram{\begin{tikzcd}[ampersand replacement=\&,row sep=small,column sep=small]
\underline{\mathrm{Ch}}_\Sigma(a,b)\times\underline{\mathrm{Ch}}_\Sigma(b,c)\arrow[r] \& \Pi^*(\mathrm{CH}_\Sigma^{\unsure{\prime}}(a,c))\arrow[d,equal]\\
\&(\underline{\mathrm{Ty}}\times\underline{\mathrm{Ty}})\underset{\underline{\mathrm{Ty}}}{\times}\underline{\mathrm{Ch}}_\Sigma(a,c)\arrow[dl,shorten <=2pt]\\
\text{produit fibré \unsure{pour|par} }\Pi:\underline{\mathrm{Ty}}\times\underline{\mathrm{Ty}}\to\underline{\mathrm{Ty}}\&
\end{tikzcd}}

Comme on a d'autre part un \deleted{\unsure{morphisme}} foncteur du dernier membre (produit fibré) vers \unsure{le} facteur $\underline{\mathrm{Ch}}_\Sigma(a,c)$, on gagne. (On l'aurait \unsure{retrouvé|retrouvée} par le raisonnement direct que je m'apprêtais à faire, et que j'ai biffé -- mais je suis plus heureux comme ça~!)

Pour bien faire, il faudrait prouver l'\uline{associativité} des foncteurs de composition, pas seulement au niveau des objets, par des arguments comme ceci, mais ça revient simplement à une relation du type
\[
(\unsure{w|u}''\mathbin{\unsure{\circ|*}}w')*w^{\unsure{4|\prime}}=w''*(w'*w)
\]
pour \deleted{\unsure{des}} un système de trois flèches de
\[
\underline{\mathrm{Ch}}_\Sigma(a,b),\quad \underline{\mathrm{Ch}}_\Sigma(b,c)\quad\text{et}\quad\underline{\mathrm{Ch}}_\Sigma(c,d)
\]
respectivement. Ça ne peut être que trivialement
